\documentclass[12pt]{article}
\usepackage[utf8]{inputenc}
\usepackage[T1]{fontenc}
\usepackage{hyperref}
\title{Zapotec, Tabaa Language Resources for K-12 Curriculum Development}
\author{Compiled by Claude Sonnet 5.5 (claude-sonnet-5-5) for Andr\'as Kornai}
\date{}
\begin{document}
\maketitle

\section{Introduction}
Tabaa Zapotec (ISO 639-3 zat; Glottocode taba1268) is an endangered Otomanguean language of Oaxaca, Mexico. The ASJP entry gives about 2,000 speakers. Searches found only partial lexical data, a manuscript grammar reference and religious audio, and no school materials.

\section{Primary Educational Resources}
No primers, textbooks or school programmes for Tabaa Zapotec were found.

\section{Linguistic and Reference Materials}
\begin{itemize}
\item ASJP wordlist for Tabaa Zapotec, which cites an unpublished Spanish-language grammar (Earl 2011) held by ILV; access to the manuscript was not confirmed. \url{https://asjp.clld.org/languages/TABAA_ZAPOTEC}
\item Wiktionary has a single Tabaa Zapotec entry, \emph{nhez} ``way, path''. \url{https://en.wiktionary.org/wiki/nhez}
\item Glottolog entry for Tabaa Zapotec. \url{https://glottolog.org/resource/languoid/id/taba1268.bigmap.html}
\item A multilingual Zapotec dictionary by \'Oscar M\'endez Espinoza, hosted in the AILLA archive, draws on many Zapotec communities; whether it covers Tabaa was not verified. \url{https://islandora-ailla.lib.utexas.edu/islandora/object/ailla%3A279600}
\end{itemize}

\section{Religious and Community Resources}
Global Recordings Network lists evangelistic and Bible-teaching audio and links to Scripture and Jesus Film resources for Tabaa Zapotec. \url{https://globalrecordings.net/en/language/zat}

\section{Digital and Audio Resources}
The Global Recordings Network page above is the only audio source found.

\section{Curriculum Development Context}
The language is endangered with few documented resources. A curriculum would require new documentation work with the community, and a first step would be to ask ILV/SIL about the Earl grammar manuscript.

\end{document}